\section{Introduction}\label{sec:intro}

An Asian option pays a function of the average price of its underlying over a fixed window. Averaging
makes the contract cheaper than its European counterpart and harder to manipulate near expiry, which is
why it is common on commodities and currencies. It also makes the contract harder to price. Even when the
underlying follows geometric Brownian motion, the arithmetic average of a lognormal process has no
tractable density, and the fixed-strike arithmetic Asian call has no closed-form price. For that model
the problem has been studied from many sides: the Laplace-transform formula of \citet{GemanYor1993} and
its extension to jump diffusions by \citet{CaiKou2012}, the spectral expansion of \citet{Linetsky2004},
and partial differential equation (PDE) methods, either on the full state space \citep{ZvanForsythVetzal1998}
or after the similarity reduction of \citet{RogersShi1995} and \citet{Vecer2001,Vecer2002}, which turns
the problem into a single space dimension. For continuous monitoring under geometric Brownian motion,
prices accurate to ten digits are available.

Under stochastic volatility the picture is thinner. In the \citet{Heston1993} model the variance is a
second state variable, the reduction of Rogers, Shi and Ve\v{c}e\v{r} leaves a two-dimensional PDE, and
the variance process may reach zero. \citet{FouqueHan2003} derived asymptotic approximations for fast
mean-reverting volatility. \citet{EwaldMenkensTing2013} obtained, through what they call the Australian
perspective, a two-dimensional pricing equation with time-independent coefficients and solved it with an
alternating direction implicit (ADI) finite difference scheme on a uniform grid, alongside Monte Carlo
schemes. \citet{FusaiKyriakou2016} developed optimised lower bounds for arithmetic Asian options in a
wide class of models and reported values for continuously monitored options under Heston. For the
geometric average under Heston, \citet{KimWee2014} gave a semi-closed-form price. For discrete monitoring
under stochastic volatility, transform and lattice-type methods are efficient
\citep{Corsaro2019,KirkbyNguyen2020}. \citet{PerottiGrzelak2024} extend Monte Carlo pricing of discretely
monitored Asian options under Heston with a data-driven sampling scheme and pay particular attention to
parameters that violate the Feller condition. \citet{CustersKarbachFriesen2026} derive semi-closed
formulas for discretely monitored Asian options in the rough Volterra--Heston model and use the geometric
price as a control variate for the arithmetic one. Recent work thus concentrates on discrete monitoring;
for the continuous case under Heston the published numbers are still mainly those of
\citet{EwaldMenkensTing2013} and \citet{FusaiKyriakou2016}. High-order discretisations of the Heston PDE exist as well, but
mainly for European and American options; a standard reference for the ADI side is
\citet{inthoutFoulon2010}.

What we did not find is a high-order solver for the continuously monitored arithmetic Asian option under
Heston whose accuracy has been checked against a reference fine enough to resolve it. The check matters
most in the case practitioners actually meet. Calibrated Heston parameters often violate the Feller
condition $2\kappa\theta\ge\xi^2$, so the variance reaches zero, the diffusion in the variance direction
degenerates on the boundary, and the solution may be less regular there. A method can converge smoothly
under refinement and still carry an error that refinement does not reveal; we meet two examples of this in
Section~\ref{sec:regimes}.

This paper offers such a solver together with the references needed to test it. The solver discretises
the reduced PDE with tensor-product quintic B-splines and collocates at Greville abscissae. Because the
coefficients of the reduced equation do not depend on time, a Crank--Nicolson march with Rannacher
start-up \citep{Rannacher1984} needs only two sparse LU factorisations, and since the strike enters only
through the point at which the solution is read, one solve prices all strikes. The boundary $v=0$, where
the Fichera theory forbids a boundary condition, is simply a line of collocation points on which the
equation degenerates by itself.

Most of the paper is about checking this solver. We use three references of increasing strength. The
first is the geometric-average version of the same contract, whose price we compute exactly (to about
$10^{-10}$) from a Riccati system; it shares the variance dynamics and the degenerate boundary with the
arithmetic problem and so tests the numerical machinery, though not the arithmetic reduction itself. The
second is the lower bound of \citet{FusaiKyriakou2016}. For Heston with a continuous average the transform
the bound requires is not given in that paper; we derive it by conditioning twice, which leads to two
consecutive Riccati systems, and validate the result independently. The third reference combines the two
previous ideas: the bound, being exact for the continuous average, serves as a control variate in a Monte
Carlo simulation, and the resulting estimator reaches a standard error of about $1.6\times10^{-5}$ with
no detectable time-discretisation bias. With this reference we compare the solver at eleven strikes and in
thirteen parameter regimes, measure its convergence in time and space, and compare it with the modified
Craig--Sneyd ADI scheme of \citet{inthoutFoulon2010} applied to the same equation.

Two findings came out of this checking and may be useful beyond this paper. First, the truncation of the
computational domain can produce errors that no refinement study will show, because refining the grid does
not move the boundary. A variance cut-off $v_{\max}=1.5$ is ample when the vol-of-vol is $\xi=1$ but
shifts the price by $7\times10^{-4}$ when $\xi=1.5$, and at a five-year maturity a cut-off $x_{\max}=8$
in the similarity variable shifts the at-the-money price by $3.5\times10^{-4}$. Only the independent reference exposed them.
Second, we could not reproduce the lower bounds that \citet{FusaiKyriakou2016} report for continuously
monitored options under Heston. Our bound is higher by up to $0.011$, a small amount in practical terms but
many standard errors of our simulation of the same quantity. We explain in Section~\ref{sec:published}
which causes we can exclude.

The paper is organised as follows. Section~\ref{sec:problem} states the problem, derives the reduced
equation and discusses its boundaries. Section~\ref{sec:method} describes the collocation scheme.
Section~\ref{sec:references} constructs the three references. Section~\ref{sec:results} reports the
numerical results, and Section~\ref{sec:conclusions} concludes. Appendix~\ref{app:transform} gives the
derivation of the transform used for the bound, and Appendix~\ref{app:tables} lists the numbers behind the
figures.
